% IEEE conference-format technical report. Compile with pdflatex + bibtex (or upload the whole report/ folder to Overleaf).
% Tables in generated/ and the result figures in figures/ are produced by scripts/make_report_assets.py from the real experiment outputs.
\documentclass[conference]{IEEEtran}
\IEEEoverridecommandlockouts
\usepackage{cite}
\usepackage{amsmath,amssymb}
\usepackage{graphicx}
\usepackage{booktabs}
\usepackage{xcolor}
\usepackage{url}
\usepackage[hidelinks]{hyperref}

\newcommand{\todo}[1]{{\color{red}\textbf{[TODO: #1]}}}
\newcommand{\gen}[1]{\IfFileExists{generated/#1.tex}{\input{generated/#1.tex}}{\textbf{[pending: \detokenize{#1}]}}}
\newcommand{\figmaybe}[2][\columnwidth]{\IfFileExists{figures/#2}{\includegraphics[width=#1]{figures/#2}}{\fbox{\parbox{0.8\dimexpr#1\relax}{\centering\small pending figure: \detokenize{#2}}}}}
\newcommand{\Ltwo}{\ensuremath{\mathcal{L}}}

\begin{document}

\title{Restoration Studio: Universal, Hard-Routed and Soft Mixture-of-Experts Image Restoration, and Style-Conditioned Face-to-Sketch Generation}

\author{\IEEEauthorblockN{Afnan Rizwan}
\IEEEauthorblockA{Generative AI, Assignment 1\\
Repository: \url{https://github.com/AfnanRizwan1/GenAI-A1}\\
Demonstration video: \todo{insert the YouTube link}}}

\maketitle

\begin{abstract}
We design, train, evaluate and deploy four related generative systems and integrate them into one web application. Three of them restore
$128\times128$ images of the Oxford-IIIT Pet dataset that are corrupted at run time by salt-and-pepper noise, Gaussian blur or rectangular
occlusion: (1) a single universal denoising autoencoder with a compressed bottleneck, (2) a corruption classifier that hard-routes each image to one of three
specialist autoencoders, and (3) a jointly trained soft mixture of experts in which a gate assigns a continuous weight to an identity branch and
to the three experts. The fourth system is (4) a style-conditioned conditional GAN that translates face photographs of the FS2K dataset into
sketches in one of three styles. Hyper-parameters of every model are tuned with Optuna, experiments are tracked with MLflow, the
inference models are exported to ONNX and verified against PyTorch, and the whole system runs from one \texttt{docker compose} command
behind a React and FastAPI application whose interface was designed in Google Stitch. \todo{add the headline numbers once the final results are in}
\end{abstract}

\begin{IEEEkeywords}
denoising autoencoder, mixture of experts, conditional GAN, image restoration, Optuna, ONNX, FastAPI
\end{IEEEkeywords}

%% ------------------------------------------------------------------------------------------------
\section{Introduction}
Real images are rarely clean: sensors add impulse noise, optics and motion blur the picture, and objects occlude parts of the scene. A
practical restoration system therefore faces a choice. One model can be trained to cope with every degradation (a universal model), or the
degradation can be recognised first and a specialist applied (hard routing). Both have drawbacks: the universal model must share
capacity among very different inverse problems, whereas hard routing fails when the classifier is wrong or when an image carries more than one degradation.
A differentiable \emph{soft} mixture of experts removes the hard decision and lets the routing be learned from the reconstruction error.

This report studies this progression on a controlled benchmark (Tasks~1--3) and adds a fourth, generative task: translating face photographs into
sketches whose style is chosen by the user (Task~4). The assignment is as much an engineering exercise as a modelling one, so the
report also documents the data pipeline, the hyper-parameter searches, the experiment tracking, the ONNX export and the deployed application.
All source code, configuration, Optuna studies, export code, Dockerfiles and the instructions to run everything are in the repository
linked on the title page.

Contributions: (i) a run-time corruption pipeline with deterministic validation and test manifests shared by all three restoration tasks;
(ii) a comparison of a universal autoencoder, oracle-routed and classifier-routed specialists and a jointly fine-tuned mixture of experts under identical conditions;
(iii) an analysis of the learned routing, including checks for inactive or dominating experts;
(iv) a style-conditioned pix2pix-style generator and discriminator in which a learned style embedding enters both networks;
(v) a containerised application that exposes all four systems and reports timing, routing and quality information.

%% ------------------------------------------------------------------------------------------------
\section{Related Work}
\textbf{Denoising autoencoders.} Vincent \emph{et al.}\ showed that reconstructing a clean input from a corrupted one forces an autoencoder to learn robust
features~\cite{vincent2008}. Our restoration models follow this idea with a convolutional encoder, a genuinely compressed code and a convolutional decoder.
We combine an $L_1$ term with a structural-similarity term~\cite{wang2004}, which rewards preserved edges and local contrast rather than only pixel agreement.

\textbf{U-Net and skip connections.} U-Net~\cite{ronneberger2015} concatenates encoder features into the decoder. This is excellent for translation tasks and is used in our
Task~4 generator, but unrestricted skip connections would let a restoration model copy the corrupted input around the bottleneck, so the Task~1--3 autoencoders have none.

\textbf{Mixtures of experts.} Jacobs \emph{et al.}\ introduced gated mixtures of local experts~\cite{jacobs1991}; Shazeer \emph{et al.}\ scaled them with a sparsely-gated layer~\cite{shazeer2017},
and Switch Transformers~\cite{fedus2022} added an auxiliary load-balancing loss to stop the router from collapsing onto few experts. Our soft mixture is dense
(every branch contributes) and uses a balance regulariser on the batch-mean routing weights in the same spirit.

\textbf{Conditional GANs and image-to-image translation.} GANs~\cite{goodfellow2014} and their conditional variants~\cite{mirza2014} underlie pix2pix~\cite{isola2017}, which pairs a
U-Net generator with a PatchGAN discriminator and an $L_1$ reconstruction term. We follow this recipe and make the condition (a sketch style) part of both networks through
a learned embedding, using feature-wise linear modulation~\cite{perez2018} in the generator.

\textbf{Hyper-parameter optimisation and tooling.} We use Optuna~\cite{akiba2019} with the tree-structured Parzen estimator~\cite{bergstra2011} and a median pruner, MLflow~\cite{mlflow} for tracking,
PyTorch~\cite{paszke2019} for modelling and ONNX Runtime~\cite{onnxruntime} for inference.

%% ------------------------------------------------------------------------------------------------
\section{Datasets and Data Pipeline}
\subsection{Oxford-IIIT Pet (Tasks 1--3)}
We use the 3\,680 images of the official \emph{trainval} split as development data and the 3\,669 images of the official \emph{test} split only for final evaluation.
The class labels are not used; the clean image is the restoration target. Development data are divided 80/20 with random seed 42 into 2\,944 training and 736 validation images, and the
same split is used by all three tasks. Every image is converted to RGB and resized to $128\times128$ once and cached.

\subsection{Run-time corruption}
During training each clean image is turned into one of four input conditions with equal probability, and a new condition and severity are drawn every time an image is loaded;
no corrupted copy is stored. The corruption is computed on the GPU on whole batches:
\begin{itemize}
\item \textbf{Salt-and-pepper:} a fraction $p\sim\mathcal{U}(0.02,0.15)$ of pixels is replaced by black or white with equal probability.
\item \textbf{Gaussian blur:} kernel size from $\{3,5,7\}$, standard deviation $\sigma\sim\mathcal{U}(0.5,2.5)$ (separable convolution with reflect padding, equal to OpenCV's \texttt{GaussianBlur}; tested to $10^{-4}$).
\item \textbf{Occlusion:} one to three black rectangles whose total area is $\mathcal{U}(10\%,35\%)$ of the image, at random positions and with aspect ratios in $[0.5,2]$.
\end{itemize}
The batched training version places rectangles independently, so overlapping rectangles make the achieved coverage slightly smaller than the target. The deterministic version used for the
manifests avoids overlap by rejection sampling and records the achieved coverage.

\subsection{Deterministic validation and test manifests}
Validation and test corruptions are generated once and stored (type, severity level, parameters, rectangle coordinates and a per-entry random seed), so every model is evaluated on exactly the same
corrupted images. The validation manifest contains, for each of the 736 validation images, the clean image and one entry per corruption type, with severity cycling through low, medium and high
inside equal thirds of the training ranges (2\,944 entries). The test manifest contains for each of the 3\,669 test images the clean image and, for every corruption, the three fixed levels required by the assignment
(36\,690 entries): salt-and-pepper $p\in\{0.03,0.08,0.15\}$; blur $(k,\sigma)\in\{(3,0.7),(5,1.5),(7,2.5)\}$; occlusion covering about $10\%$, $20\%$ and $35\%$ with one, two and three rectangles.
Replaying a manifest is deterministic, and regenerating it reproduces the committed files byte for byte.

\subsection{FS2K (Task 4)}
FS2K~\cite{fan2022} contains 2\,104 photograph--sketch pairs in three sketch styles. We use the official training portion (1\,058 pairs) and test portion (1\,046 pairs). Fifteen percent of the training pairs, stratified by style with seed 42,
form the validation set, which gives 899 training pairs and 159 validation pairs (54, 52 and 53 per style). The test set is style-skewed (619, 381 and 46 images for styles~1, 2 and~3) and is never used for training or tuning.
Photographs and sketches are resized to $128\times128$ and the pairing is preserved by index. Augmentation (horizontal flip, rotation up to $5^\circ$, zoom $0.9$--$1.0$, shift up to $5\%$) is applied
by a single affine warp to the six-channel stack of photograph and sketch, so the pixel correspondence is preserved by construction.

%% ------------------------------------------------------------------------------------------------
\section{Task 1: Universal Multi-Corruption Denoising Autoencoder}
\subsection{Architecture}
The encoder has four stages, each a stride-2 $4\times4$ convolution followed by a $3\times3$ convolution (batch normalisation~\cite{ioffe2015} and LeakyReLU after each), taking the spatial size from $128$ to $8$ while the channels
grow from $c$ to $8c$. The bottleneck is one of two kinds, chosen by Optuna: a \emph{linear} code (flatten, then a fully connected layer to $d$ numbers) or a \emph{convolutional} code (a $1\times1$ convolution to $c_z$ channels, i.e.\ a
$c_z\times8\times8$ tensor of $64c_z$ numbers). Dropout is applied to the code. The decoder mirrors the encoder with transposed convolutions and ends with a sigmoid. Even the smallest code is more than ten times smaller than the $49\,152$ input values.
There are \emph{no skip connections}, so every output pixel depends only on the code (a unit test checks that decoding the encoder output reproduces the forward pass). Because no skip connection is used there is
nothing to justify or ablate in the sense of the assignment's remark on limited skips; a skip-connection ablation was not run.

\subsection{Loss, optimiser and tuning}
With $x$ the clean image, $\tilde{x}$ the corrupted input and $\hat{x}=D(E(\tilde{x}))$ the reconstruction,
\begin{equation}
\Ltwo=\alpha\,\lVert\hat{x}-x\rVert_1+(1-\alpha)\bigl(1-\mathrm{SSIM}(\hat{x},x)\bigr),
\end{equation}
with SSIM computed over $11\times11$ Gaussian windows ($\sigma=1.5$). Models are trained with AdamW~\cite{loshchilov2019}, a one-cycle learning-rate schedule~\cite{smith2019}, mixed precision, and the whole dataset resident on the GPU.
The Optuna study (TPE sampler, seed 42, median pruner) searches the learning rate $[3\cdot10^{-4},3\cdot10^{-3}]$ (log), batch size $\{32,64,128\}$, base channels $c\in\{16,32,48,64\}$, dropout $[0,0.3]$, the loss weight $\alpha\in[0.5,0.95]$, the bottleneck type,
and its size ($d\in\{128,256,512,1024\}$ or $c_z\in\{4,8,16,32\}$). The objective, minimised on the validation manifest, is $L_1+(1-\mathrm{SSIM})$, which is independent of $\alpha$ so that trials with different $\alpha$ are comparable.
The planned budget is 30 trials of 10 epochs followed by 150 epochs of final training. \todo{state the completed/pruned trial counts, best trial and final configuration from Table~\ref{tab:optuna}}

\subsection{A first run and the bottleneck study}
A first, smaller run (8 trials of 5 epochs, linear bottleneck only, 40 final epochs) reached a validation PSNR of 18.75\,dB and SSIM of 0.471, and the curve had flattened. Because the linear code discards spatial layout, we added the convolutional bottleneck to the
search space for the main run. \todo{compare the main run with this first run: PSNR/SSIM gain, which bottleneck Optuna selected}

\subsection{Results}
Tables~\ref{tab:rest_psnr} and~\ref{tab:rest_ssim} report PSNR and SSIM on the fixed test manifest for every corruption and severity, together with the no-restoration baseline (the corrupted input itself), which every model must beat.
\figmaybe[\textwidth]{curves_t1_t2.pdf}
\todo{interpret the Task 1 rows: where the autoencoder beats the input and where it does not, per corruption and severity; discuss the strict-bottleneck limitation}

\begin{figure*}[t]
\centering
\figmaybe[0.82\textwidth]{examples_t1.png}
\caption{Twelve test examples (clean target, corrupted input, Task~1 reconstruction, absolute error map) covering every corruption and severity, plus extra high-severity cases.}
\label{fig:t1_examples}
\end{figure*}
\begin{figure}[t]
\centering
\figmaybe{failures_t1.png}
\caption{Task~1 failure cases: for each corruption the test entry with the lowest PSNR gain over the input, and the worst remaining entry.}
\label{fig:t1_failures}
\end{figure}
\todo{discuss the four failure cases in Fig.~\ref{fig:t1_failures}: what went wrong and why}

%% ------------------------------------------------------------------------------------------------
\section{Task 2: Corruption Classification and Hard-Routed Specialists}
\subsection{Classifier}
The classifier has four blocks of a stride-1 and a stride-2 $3\times3$ convolution (batch normalisation, ReLU), global average pooling, dropout and a linear layer with four outputs. Its labels come from the run-time corruption pipeline, and every training batch is exactly balanced
(equal numbers of clean, salt-and-pepper, blur and occlusion samples), so the classifier cannot favour one class. It is trained with cross-entropy and AdamW. The Optuna study searches the learning rate, batch size, channel configuration
(small, medium, large), dropout and weight decay, and maximises validation macro-F1. \todo{report the study (Table~\ref{tab:optuna}) and the final configuration}

\subsection{Specialists and hard routing}
Three specialist autoencoders share one architecture but are trained independently, each only on its own corruption and with the same clean targets. To keep the cost feasible the architecture is chosen by one shared Optuna search on the three corruptions
(clean images excluded) with the same space and objective as Task~1; the three specialists are then trained from scratch with that configuration. At inference the classifier's arg-max selects the specialist, and a clean prediction uses an identity bypass
(no expert is called).

\subsection{Evaluation modes}
\emph{Oracle routing} selects the expert from the true label in the test manifest and measures the specialists in isolation. \emph{Predicted routing} uses the classifier and measures the operational system. The difference between the two is the cost of classifier errors.

\subsection{Results}
Test accuracy of the classifier is \gen{classifier_summary} (Table~\ref{tab:cls}, Fig.~\ref{fig:confusion}); \gen{misrouted} of the test entries were misrouted.
\begin{table}[t]
\caption{Task~2 classifier on the test manifest (per class, macro average).}
\label{tab:cls}
\centering\small
\gen{classifier_metrics}
\end{table}
\begin{table}[t]
\caption{Classifier accuracy by corruption and severity (\%).}
\label{tab:cls_cond}
\centering\small
\gen{classifier_by_condition}
\end{table}
\begin{figure}[t]
\centering
\figmaybe[0.8\columnwidth]{confusion_matrix.png}
\caption{Row-normalised confusion matrix of the four-class classifier on the test manifest.}
\label{fig:confusion}
\end{figure}
\todo{interpret the per-class metrics, the confusion matrix and the accuracy-by-severity table (which confusions occur and at which severity)}

\begin{figure}[t]
\centering
\figmaybe{failures_t2_misrouting.png}
\caption{Misrouting cases: the four classifier errors with the largest PSNR loss compared with oracle routing.}
\label{fig:misroute}
\end{figure}
\todo{discuss the misrouting cases in Fig.~\ref{fig:misroute} and the gap between oracle and predicted routing in Tables~\ref{tab:rest_psnr}--\ref{tab:rest_ssim}}

%% ------------------------------------------------------------------------------------------------
\section{Task 3: Jointly Trained Soft Mixture of Experts}
\subsection{Model}
Let $z(\tilde{x})\in\mathbb{R}^4$ be the logits of the gate and $E_s,E_b,E_o$ the salt-and-pepper, blur and occlusion experts. With temperature $T$,
\begin{align}
g &= \mathrm{softmax}\bigl(z(\tilde{x})/T\bigr)=[g_0,g_s,g_b,g_o],\\
\hat{x} &= g_0\,\tilde{x}+g_s E_s(\tilde{x})+g_b E_b(\tilde{x})+g_o E_o(\tilde{x}).
\end{align}
The first branch is an identity branch for clean inputs. The combination is differentiable, so gate and experts are trained through the final reconstruction error. The gate is initialised from the Task~2 classifier and the experts from the three trained Task~2 specialists; nothing starts from random weights.

\subsection{Training}
Training has two stages: a warm-up in which the experts are frozen (and kept in evaluation mode, so their batch-norm statistics do not change) and only the gate is trained, followed by joint fine-tuning of all components with a smaller learning rate. The loss is
\begin{equation}
\Ltwo = a\,\Ltwo_1+(1-a)\bigl(1-\mathrm{SSIM}\bigr)+\gamma\,\mathrm{CE}\!\left(\tfrac{z}{T},c\right)+\delta\,\Ltwo_{\mathrm{bal}},
\end{equation}
where $c$ is the true corruption class and
\begin{equation}
\Ltwo_{\mathrm{bal}}=\sum_{i=1}^{K}\left(\bar{g}_i-\tfrac{1}{K}\right)^2,\qquad K=4,
\end{equation}
with $\bar{g}_i$ the mean weight of branch $i$ over a class-balanced batch, so the target for a healthy gate is uniform usage of the four branches over the batch. The assignment's starting weights are $a=0.8$ ($\beta=0.2$ for SSIM, here $1-a$), $\gamma=0.1$ and $\delta=0.01$.
(The assignment text gives the balance formula only in a garbled form, so we use the squared deviation from uniform usage, the simplest differentiable regulariser of this kind; Switch Transformers~\cite{fedus2022} use a related load-balancing loss.)
Optuna tunes the joint learning rate $[10^{-5},5\cdot10^{-4}]$ (log), temperature $T\in[0.5,3]$, classification weight $\gamma\in[0.01,1]$ (log), balance weight $\delta\in[10^{-3},1]$ (log) and $a\in[0.5,0.95]$.
A trial is pruned when its validation objective is poor, or when one branch receives a mean weight above 0.9 after the warm-up (routing collapse). \todo{state the completed / pruned trial counts and the selected configuration}

\subsection{Routing analysis}
\gen{moe_gate_acc} of test entries have their largest weight on the branch that matches the true condition. Table~\ref{tab:moe_w} lists the mean routing weights per true corruption and severity and Fig.~\ref{fig:heatmap} shows them as a heat map.
Table~\ref{tab:moe_act} checks for an inactive expert (mean weight below 0.05 over all inputs) and for an expert that dominates unrelated inputs (mean weight above 0.5 on inputs of other corruption types).
\begin{table}[t]
\caption{Task~3: mean routing weights on the test manifest by true corruption and severity (largest per row in bold).}
\label{tab:moe_w}
\centering\small
\gen{moe_weights}
\end{table}
\begin{table}[t]
\caption{Task~3: expert activity checks over all test entries.}
\label{tab:moe_act}
\centering\scriptsize
\gen{moe_activity}
\end{table}
\begin{figure}[t]
\centering
\figmaybe[0.9\columnwidth]{moe_routing_heatmap.png}
\caption{Mean routing weights by true corruption and severity.}
\label{fig:heatmap}
\end{figure}
\begin{figure*}[t]
\centering
\begin{minipage}{0.48\textwidth}\centering\figmaybe{moe_examples_dominant.png}\end{minipage}\hfill
\begin{minipage}{0.48\textwidth}\centering\figmaybe{moe_examples_distributed.png}\end{minipage}
\caption{Left: inputs for which one expert dominates (lowest routing entropy). Right: inputs for which the weights are spread over several branches (highest entropy). Bars show the four routing weights.}
\label{fig:moe_examples}
\end{figure*}
\todo{interpret Tables~\ref{tab:moe_w}--\ref{tab:moe_act} and Figs.~\ref{fig:heatmap}--\ref{fig:moe_examples}: which expert handles which corruption, whether any expert is inactive or dominating, and how the soft system compares with hard routing in Tables~\ref{tab:rest_psnr}--\ref{tab:rest_ssim}}

%% ------------------------------------------------------------------------------------------------
\begin{table*}[t]
\caption{Restoration quality on the fixed test manifest: PSNR (dB). The first column is the corrupted input without restoration. Best per row in bold (clean images give 60\,dB for an identity mapping, the cap used for identical images).}
\label{tab:rest_psnr}
\centering\small
\resizebox{0.78\textwidth}{!}{\gen{restoration_psnr}}
\end{table*}
\begin{table*}[t]
\caption{Restoration quality on the fixed test manifest: SSIM. Best per row in bold.}
\label{tab:rest_ssim}
\centering\small
\resizebox{0.78\textwidth}{!}{\gen{restoration_ssim}}
\end{table*}

%% ------------------------------------------------------------------------------------------------
\section{Task 4: Style-Conditioned Face-to-Sketch Generation}
\subsection{Architecture}
\textbf{Generator.} A U-Net with six down-sampling and six up-sampling levels ($128\to2\to128$) and skip connections. The style $s\in\{1,2,3\}$ selects a learned embedding $e_s\in\mathbb{R}^{m}$ ($m$ tuned). The embedding enters the generator in two places, so it is part of the model and not an interface label:
it is broadcast over the image and concatenated to the input photograph, and in every decoder block it produces a per-channel scale and shift (FiLM~\cite{perez2018}) applied after instance normalisation~\cite{ulyanov2016}.
\textbf{Discriminator.} A PatchGAN that receives the photograph, the (real or generated) sketch and the broadcast style embedding and outputs a grid of real/fake logits ($14\times14$ for $128\times128$ inputs). A unit test confirms that the style embedding receives gradients in both networks.

\subsection{Objective and training}
With photograph $x$, style $s$ and target sketch $y$, the discriminator minimises the average of the binary cross-entropy on real pairs $(x,s,y)$ and fake pairs $(x,s,G(x,s))$ (logged separately as D-real and D-fake) and the generator minimises
\begin{equation}
\Ltwo_G=\mathrm{BCE}\bigl(D(x,s,G(x,s)),1\bigr)+\lambda\,\lVert G(x,s)-y\rVert_1 ,
\end{equation}
with adversarial and reconstruction parts logged separately. Both networks use Adam ($\beta=(0.5,0.999)$). Optuna searches the generator and discriminator learning rates $[5\cdot10^{-5},10^{-3}]$ (log), batch size $\{8,16,32\}$, base channels $\{16,32,48,64\}$, dropout $[0,0.5]$,
embedding dimension $\{8,16,32,64\}$ and $\lambda\in[10,200]$ (log), using 15 epochs per trial; the best configuration is retrained for 150 epochs. The objective is validation $L_1+(1-\mathrm{SSIM})$.
The same first eight validation photographs are rendered every ten epochs to follow the generator's development.

\subsection{Results}
Table~\ref{tab:t4} reports $L_1$, SSIM and PSNR on the official test split per style, together with a baseline that simply converts the photograph to grayscale. Table~\ref{tab:t4_cond} checks style conditioning: every test photograph is generated with all three styles and compared with its ground-truth sketch.
\begin{table}[t]
\caption{Task~4 on the official FS2K test split (generator vs.\ grayscale-photo baseline, subscript $g$).}
\label{tab:t4}
\centering\scriptsize
\gen{t4_metrics}
\end{table}
\begin{table}[t]
\caption{Style conditioning: mean $L_1$ to the ground-truth sketch (rows: true style, columns: style given to the generator; lowest per row in bold).}
\label{tab:t4_cond}
\centering\small
\gen{t4_conditioning}
\end{table}
\begin{figure*}[t]
\centering
\figmaybe[0.7\textwidth]{t4_examples_test.png}
\caption{Test photographs generated in each of the three styles next to the ground-truth sketch of their true style.}
\label{fig:t4_examples}
\end{figure*}
\begin{figure}[t]
\centering
\figmaybe[0.85\columnwidth]{t4_failures_test.png}
\caption{Task~4 failure cases: the four test pairs with the largest $L_1$ error.}
\label{fig:t4_fail}
\end{figure}
\todo{interpret Tables~\ref{tab:t4}--\ref{tab:t4_cond}, the loss curves (Fig.~\ref{fig:curves34}) and the example and failure grids; note the small style-3 test set (46 images) and the limits of an $L_1$-dominated objective on 899 training pairs}

\begin{figure*}[t]
\centering
\figmaybe[0.95\textwidth]{curves_t3_t4.pdf}
\caption{Left: Task~3 validation PSNR (dashed line: start of joint fine-tuning). Middle: gate accuracy and the largest mean routing weight (collapse indicator). Right: Task~4 training losses (D-real, D-fake, G-adversarial, G-reconstruction).}
\label{fig:curves34}
\end{figure*}

%% ------------------------------------------------------------------------------------------------
\section{Hyper-parameter Optimisation}
Every task uses an Optuna study (SQLite storage, TPE sampler with seed 42, median pruner, resumable) whose trials, parameters and intermediate values are also logged to MLflow. Table~\ref{tab:optuna} summarises the studies and lists the selected parameters.
\begin{table}[t]
\caption{Optuna studies: number of trials and best objective value (validation $L_1+(1-\mathrm{SSIM})$ for the autoencoders and the GAN, macro-F1 for the classifier; lower is better except for the classifier).}
\label{tab:optuna}
\centering\small
\gen{optuna_summary}
\end{table}
\begin{itemize}\footnotesize
\IfFileExists{generated/optuna_best.tex}{\item \textbf{Task 1 autoencoder}: lr=0.000774, batch\_size=32, base\_ch=48, dropout=0.1006, alpha=0.5137, bottleneck=conv, latent\_ch=32
\item \textbf{Task 2 classifier}: lr=0.003663, batch\_size=32, channels=small, dropout=0.4426, weight\_decay=9.517e-05
\item \textbf{Task 2 specialists (shared)}: lr=0.001367, batch\_size=32, base\_ch=48, dropout=0.04377, alpha=0.6252, bottleneck=conv, latent\_ch=32
\item \textbf{Task 3 soft MoE}: lr\_joint=0.0002043, temperature=2.175, gamma=0.01208, delta=0.0498, a=0.5275
\item \textbf{Task 4 cGAN}: lr\_g=0.0009958, lr\_d=0.0003826, batch\_size=8, base\_ch=48, dropout=0.1291, emb\_dim=8, lam=196.2
}{\item \textbf{[pending: Optuna results]}}
\end{itemize}

%% ------------------------------------------------------------------------------------------------
\section{Application and Deployment}
\subsection{Design}
The interface was designed in Google Stitch~\cite{stitch} before it was implemented (project: \url{https://stitch.withgoogle.com/projects/10012210425231381214}). The first generation (Fig.~\ref{fig:stitch1}) already fixed the layout (sidebar with four workspaces, controls card, input / output panels, metric cards), but it
contained details the system does not have (hardware names, network architectures, extra sketch controls, scores). We therefore wrote a design system (\texttt{design/DESIGN.md}: palette, typography, spacing, components, motion and strict rules about which data may be shown) and generated a refined design from it
(Fig.~\ref{fig:stitch2}). The React and Tailwind implementation follows the second design and shows only quantities the application really computes.
\begin{figure}[t]
\centering
\figmaybe{stitch_v1_universal.png}
\caption{Original Stitch design (first generation), Universal Restoration screen.}
\label{fig:stitch1}
\end{figure}
\begin{figure*}[t]
\centering
\begin{minipage}{0.24\textwidth}\centering\figmaybe{stitch_v2_universal.png}\end{minipage}\hfill
\begin{minipage}{0.24\textwidth}\centering\figmaybe{stitch_v2_hard.png}\end{minipage}\hfill
\begin{minipage}{0.24\textwidth}\centering\figmaybe{stitch_v2_soft.png}\end{minipage}\hfill
\begin{minipage}{0.24\textwidth}\centering\figmaybe{stitch_v2_sketch.png}\end{minipage}
\caption{Refined Stitch design: Universal Restoration, Hard-Routed Restoration, Soft Mixture-of-Experts and Face-to-Sketch Generator.}
\label{fig:stitch2}
\end{figure*}

\subsection{Architecture}
The frontend (React, Tailwind CSS, served by nginx) talks to a FastAPI backend that exposes \texttt{GET /api/health}, \texttt{POST /api/universal}, \texttt{/api/hard}, \texttt{/api/soft} and \texttt{/api/sketch}. The backend validates every upload (content type, size, pixel count, decodability, minimum size),
resizes it to $128\times128$, optionally applies a chosen corruption (type, severity or custom parameters, seed) using exactly the corruption code that was used for training, runs the ONNX models with ONNX Runtime on the CPU and returns the images together with the corruption settings,
inference time, PSNR and SSIM against the clean reference, and the classifier probabilities or routing weights (and their entropy). Already corrupted images can be uploaded and restored as they are. A difference map, a before/after slider and JSON export of the metrics are available.
\figmaybe[\textwidth]{app_screenshots.png}
\todo{insert screenshots of the running application with the trained models (one per workspace) and replace this placeholder}

\subsection{ONNX export and verification}
All inference models (universal autoencoder, classifier with softmax, the three specialists, the complete soft mixture as one graph returning the image and the four weights, and only the generator of Task~4) are exported with opset 17 and a dynamic batch axis. Each exported model is run with ONNX Runtime and compared with PyTorch on several inputs and two batch sizes;
the export command fails if the maximum absolute difference exceeds $10^{-4}$ (Table~\ref{tab:onnx}).
\begin{table}[t]
\caption{ONNX export verification (ONNX Runtime vs.\ PyTorch).}
\label{tab:onnx}
\centering\scriptsize
\resizebox{\columnwidth}{!}{\gen{onnx}}
\end{table}

\subsection{Containers, tracking and reproducibility}
\texttt{docker compose up --build} builds a backend image (Python, ONNX Runtime, no PyTorch) and a frontend image (multi-stage Node build served by nginx, which also proxies \texttt{/api}); the trained models are mounted from \texttt{./models}. The Compose setup was verified from a fresh clone on Windows with Docker Engine running in WSL~2.
All runs were logged to MLflow (parameters, losses, metrics, checkpoints and sample images). The repository contains 59 Python tests and 33 frontend tests covering the corruption pipeline, models, training loops, Optuna objectives, evaluation, ONNX export, the API and the user interface.
Trained model files are larger than GitHub's file limit and are distributed as downloadable assets (\texttt{scripts/download\_models.sh}).

%% ------------------------------------------------------------------------------------------------
\section{Limitations}
The hyper-parameter searches are small (tens of trials of a few epochs) because of the time available, so the selected configurations are good rather than optimal; for the Task~4 search the best trial lay at the upper end of the allowed range for $\lambda$.
All models work at $128\times128$. The strict bottleneck of the restoration autoencoders limits reconstruction detail, so for mild corruptions the corrupted input itself can be closer to the clean image than the restored one (see the baseline column in Tables~\ref{tab:rest_psnr}--\ref{tab:rest_ssim}).
Training corruptions place occlusion rectangles independently, whereas the test manifest avoids overlap. The mixture of experts handles combined corruptions only to the extent its gate generalises; we did not construct a benchmark of images with several simultaneous corruptions.
The sketch generator is trained on 899 pairs with an $L_1$-dominated objective and the style-3 test set has only 46 images, so per-style numbers must be read with care. No skip-connection ablation and no repeated runs with different seeds were performed, so no confidence intervals are given.

\section{Conclusion}
\todo{write the conclusion from the final results: which restoration strategy worked best per corruption, what the routing analysis showed, how well the sketch generator uses its style condition, and what the application delivers}

%% ------------------------------------------------------------------------------------------------
\appendices
\section{Use of AI tools}
\begin{itemize}
\item \textbf{Claude Code (Anthropic, \cite{claudecode})} was used throughout as a coding assistant: to write and refactor the data pipeline, models, training and evaluation scripts, Optuna studies, ONNX export, the FastAPI backend, the React frontend, the Dockerfiles, the tests, the report skeleton and the figure and table generation scripts; and to discuss design decisions (for example the convolutional bottleneck, the balance regulariser, the FiLM conditioning).
\item \textbf{Google Stitch~\cite{stitch}} generated the interface mock-ups from prompts and a design-system file written for the project; the generated designs were corrected by hand where they contained invented content.
\end{itemize}
\textbf{How the outputs were tested and corrected.} Generated code was run, not trusted: the repository has 59 Python and 33 frontend tests, including tests that compare the batched GPU blur with OpenCV, check that decoding the code reproduces the autoencoder output (no skip path), verify that experts are frozen during warm-up, check that the augmentation transforms photograph and sketch identically, compare ONNX with PyTorch outputs, and check that the backend SSIM equals the library definition used in training.
Several failures found this way were fixed (for example a wrongly ignored \texttt{src/data} folder that broke the cloud runs, a metric key naming slip in the routing analysis, and layout problems found by driving the application in a real browser), and a mutation check confirmed that the frontend tests fail when the code is deliberately broken.
Training was run by the author on Kaggle GPUs, and every result in this report is read from the files written by the evaluation scripts. The student remains responsible for the submitted system. \todo{review and edit this appendix so it describes your own contribution accurately}

\section{Reproducing the results}
\begin{enumerate}\footnotesize
\item \texttt{bash scripts/run\_t1\_t2.sh} (Task~1 search and training, classifier, specialists), \texttt{bash scripts/run\_t3.sh}, \texttt{bash scripts/run\_t4.sh}.
\item \texttt{python -m src.eval.evaluate ...} and \texttt{python -m src.eval.evaluate\_t4 ...} (test-set evaluation), \texttt{python -m src.export.export\_onnx ...} (ONNX).
\item \texttt{python scripts/make\_report\_assets.py} regenerates every table and figure of this report from the outputs.
\item \texttt{docker compose up --build} starts the application at \url{http://localhost:8080}.
\end{enumerate}

\bibliographystyle{IEEEtran}
\bibliography{references}

\end{document}
